(1) 顶角为 \(A\in(0,\pi)\) 的均匀透明棱镜置于空气中，折射率为 \(n>1\)，光线从第一面入射、从第二面出射。定义入射角 \(i\)、出射角 \(e\)、内部相对法线角 \(r_1,r_2\)，推导总偏向角 \(\delta\) 的几何关系。
(2) 在可行且两侧可传播的对称最小偏向路径下，证明 \(r_1=r_2=A/2\)，推导最小偏向角 \(\delta_{\min}\) 与折射率关系。
(3) 对固定棱镜顶角求角分散导数 \(d\delta_{\min}/dn\)，进而求两个相近波长 \(\lambda,\lambda+d\lambda\) 的一阶色散偏向差。
(4) 说明 \(n\sin(A/2)<1\) 的几何可行条件，检查 \(n\to1\) 的零偏向极限与顶角很小时的薄棱镜近似。
